\chapter{视觉定位与多传感器融合}

\section{视觉处理总体流程}

当前视觉代码由业务层、项目接口层和底层几何层组成，覆盖摄像头采集、二值化、连通域、
信标与车标识别、鱼眼射线查表、IPM、双相机融合和世界坐标输出。

\begin{figure}[htbp]
  \centering
  \begin{tikzpicture}[node distance=7mm,box/.style={draw,rounded corners,align=center,minimum width=33mm,minimum height=9mm},line/.style={-{Latex},thick}]
    \node[box] (raw) {188×120灰度图};
    \node[box,right=of raw] (pre) {标定/阈值处理};
    \node[box,right=of pre] (blob) {连通域与候选筛选};
    \node[box,below=of blob] (marker) {信标/车标几何评分};
    \node[box,left=of marker] (ray) {鱼眼射线与IPM};
    \node[box,left=of ray] (fusion) {双相机融合与世界坐标};
    \draw[line] (raw)--(pre);
    \draw[line] (pre)--(blob);
    \draw[line] (blob)--(marker);
    \draw[line] (marker)--(ray);
    \draw[line] (ray)--(fusion);
  \end{tikzpicture}
  \caption{视觉定位处理流程}
\end{figure}

\section{摄像头标定}

仓库现有标定使用OpenCV鱼眼模型。33幅可读图像中16幅进入求解，当前RMS重投影误差为0.222251像素。

\begin{table}[htbp]
  \centering
  \caption{当前摄像头标定参数（定稿前需确认与上场相机一致）}
  \begin{tabular}{cccc}
    \toprule
    参数 & 数值 & 参数 & 数值 \\
    \midrule
    $f_x$ & 79.393980 & $f_y$ & 79.572640 \\
    $c_x$ & 95.407813 & $c_y$ & 59.954799 \\
    $k_1$ & -0.00872509 & $k_2$ & -0.02498306 \\
    $k_3$ & 0.02068140 & $k_4$ & -0.00718848 \\
    \bottomrule
  \end{tabular}
\end{table}

\todoitem{补充标定板规格、方格尺寸、采集距离/姿态范围、每视图误差分布、独立验证集误差和标定日期。}

\section{信标与车标检测}

说明动态Blob门限、几何精筛、臂长与对称性、开口比、凸包凹陷、方向先验和综合评分。
对应实现位于\codepath{beacon_detect.c}、\codepath{vehicle_marker_localization.c}和
\codepath{beacon_adaptive_threshold.c}。

\section{IPM与坐标变换}

像素通过鱼眼模型转换为相机射线，再结合相机外参、无人机姿态和高度求取地面交点。
正文应明确世界系、机体系、相机系和车体系的轴向、旋转顺序及单位。

\section{双相机融合与失效处理}

说明双相机视场选择、重叠区域融合、时间同步、跨核传输、置信度以及单相机失效时的降级策略。

\section{光流、ToF与位置观测融合}

当前默认光流为UP-FLOW-302，更新周期25 ms（40 Hz），PMW3901作为兼容路径保留。
位置环由视觉/红外定位输出目标速度，再由光流速度反馈闭环。应给出观测超时、尺度补偿、旋转补偿和滤波延迟。

\section{视觉性能实验}

\begin{table}[htbp]
  \centering
  \caption{视觉实验计划与待填写结果}
  \begin{tabularx}{\textwidth}{p{3.1cm}p{3.3cm}p{3.1cm}X}
    \toprule
    指标 & 测试条件 & 结果 & 统计方式 \\
    \midrule
    信标检出率/误检率 & 光照、距离、姿态分组 & 待测试 & 每组样本数与混淆矩阵 \\
    车标位置误差 & 标定点或动捕真值 & 待测试 & RMSE/最大误差/95分位 \\
    车标航向误差 & 多角度静态与动态 & 待测试 & 平均绝对误差/最大误差 \\
    单帧处理耗时 & CM7\_1实际运行 & 待测试 & 均值/最大值/超期次数 \\
    双相机融合连续性 & 跨视场运动 & 待测试 & 跳变量与丢失时长 \\
    \bottomrule
  \end{tabularx}
\end{table}
